\documentclass[11pt,a4paper]{article}

\usepackage[utf8]{inputenc}
\usepackage[T1]{fontenc}
\usepackage[brazilian]{babel}
\usepackage[margin=2cm]{geometry}
\usepackage{enumitem}
\usepackage{hyperref}
\usepackage{titlesec}
\usepackage{parskip}

\hypersetup{
  colorlinks=true,
  linkcolor=black,
  urlcolor=black,
  pdftitle={Kamille Hartmann Maccagnan - Currículo},
  pdfauthor={Kamille Hartmann Maccagnan}
}

\titleformat{\section}{\large\bfseries}{}{0em}{}[\titlerule]
\titlespacing*{\section}{0pt}{12pt}{6pt}

\pagestyle{empty}
\setlength{\parindent}{0pt}
\setlength{\parskip}{6pt}

\newcommand{\cvheader}[4]{%
  \begin{center}
    {\LARGE\bfseries #1}\\[4pt]
    {\large #2}\\[6pt]
    #3\\[2pt]
    #4
  \end{center}
  \vspace{4pt}
}

\newcommand{\experience}[3]{%
  \textbf{#1} \hfill \textit{#2}\\[4pt]
  #3
  \vspace{8pt}
}

\begin{document}

\cvheader{Kamille Hartmann Maccagnan}%
{Analista de Negócios | Documentação de Requisitos | Metodologias Ágeis | Gestão de Demandas}%
{Campo Comprido, Curitiba-PR \textbullet{} (41) 9 9928-4424 \textbullet{} \href{mailto:kamillehartmannm@gmail.com}{kamillehartmannm@gmail.com}}%
{\href{https://linkedin.com/in/kamille-hartmann-maccagnan/}{linkedin.com/in/kamille-hartmann-maccagnan}}

\section*{Resumo Profissional}

Profissional com experiência na interface entre negócio, operações e tecnologia, atuando no levantamento e análise de demandas, documentação funcional, mapeamento e padronização de processos e apoio à tomada de decisão. Vivência em ambientes corporativos e multinacionais, com participação em PMO, governança de dados e acompanhamento de portfólio de projetos, incluindo controle de cronogramas, indicadores de desempenho e elaboração de relatórios executivos. Experiência na tradução de necessidades operacionais em documentação clara para times técnicos, com elaboração de POPs, manuais e materiais de apoio, além de condução de treinamentos e capacitação de usuários. Atuação na organização e classificação de demandas, identificação de riscos e oportunidades de melhoria, monitoramento de métricas e facilitação de alinhamentos entre stakeholders de diferentes áreas. Conhecimento em metodologias ágeis (Scrum e Kanban) e ferramentas de gestão de demandas e projetos, com perfil analítico, organizado e comunicativo.

\section*{Experiência Profissional}

\experience{Anjuss -- Analista de Suporte de TI}{07/2025 -- Atual}{%
Atuação como ponte entre áreas de negócio, lojas e tecnologia, com foco na organização, análise e priorização de demandas internas e no suporte funcional a usuários e sistemas. Responsável pela padronização de processos e pela elaboração de documentação funcional, incluindo POPs do sistema para as lojas e POPs técnicos para a área de T.I., garantindo rastreabilidade entre necessidades operacionais, regras de uso e orientações para as equipes. Liderança do treinamento e onboarding de novos funcionários, com criação de materiais de apoio para disseminação de conhecimento e aderência aos processos definidos. Desenvolvimento de dashboards em Power BI e relatórios gerenciais para monitoramento de indicadores operacionais e acompanhamento de entregas, contribuindo para maior visibilidade das demandas e apoio à gestão. Atuação na identificação de oportunidades de melhoria de processos e na automação de rotinas, com foco em eficiência operacional e alinhamento entre stakeholders.}

\experience{HORSE do Brasil -- Estagiária de TI (IS/IT LATAM)}{07/2024 -- 07/2025}{%
Atuação em múltiplas frentes corporativas, incluindo PMO, governança de dados, suporte e interface com áreas operacionais ligadas ao chão de fábrica, em ambiente multinacional e de projetos. Acompanhamento de portfólio de projetos com organização de cronogramas, indicadores de desempenho e visibilidade de entregas, com elaboração de relatórios executivos e materiais de acompanhamento para apoio à tomada de decisão. Atuação como interface entre áreas de negócio, fornecedores e stakeholders internacionais, participando de reuniões de alinhamento conduzidas em até três idiomas simultaneamente e contribuindo para validação de entregas, classificação de demandas e compartilhamento de análises da operação. Participação em iniciativas de melhoria contínua, análise de riscos, documentação de processos e evolução de práticas de gestão e governança, com monitoramento de métricas relacionadas à evolução de projetos e iniciativas estratégicas.}

\experience{Restaurante Familiar -- Analista de Dados e Suporte Técnico}{01/2023 -- 07/2024}{%
Atuação na análise de necessidades de gestão do negócio, com levantamento de informações sobre vendas, performance operacional e resultados financeiros para apoio à tomada de decisão. Mapeamento e estruturação de processos de coleta e organização de dados, incluindo a criação da base de dados do zero em Excel para viabilizar dashboards em Power BI e definição de indicadores de desempenho. Responsável pela padronização de informações e identificação de inconsistências que impactavam a confiabilidade das análises, com proposição de melhorias em processos operacionais. Atuação no alinhamento entre áreas operacionais e administrativas, traduzindo demandas do negócio em estruturas de dados e relatórios que apoiavam a gestão e a evolução dos resultados.}

\section*{Projetos e Conquistas}

Estruturação de KPIs e dashboards executivos para monitoramento de indicadores de desempenho e suporte à tomada de decisão. Participação em iniciativas de melhoria contínua e otimização de processos com foco em eficiência operacional. Automação e padronização de fluxos de dados e relatórios, reduzindo retrabalho e aumentando a confiabilidade das informações. 1º lugar no Transformation Day Renault 2023 com solução orientada a dados e melhoria de processos. Participação no Hackathon Bosch 2023 com foco em resolução de problemas de negócio com uso de dados.

\section*{Habilidades Técnicas}

\textbf{Análise de Negócios e Requisitos:} Levantamento e análise de demandas, documentação funcional, elaboração de POPs, manuais e materiais de apoio, mapeamento e modelagem de processos, classificação de demandas, validação funcional de entregas, suporte à tomada de decisão.

\textbf{Metodologias Ágeis e Gestão de Demandas:} Scrum, Kanban, Boas Práticas PMOBOK, organização e priorização de demandas, acompanhamento de entregas, melhoria contínua, gestão de portfólio de projetos (PMO).

\textbf{Indicadores, Métricas e Reportes:} KPIs, dashboards executivos, relatórios gerenciais e de acompanhamento, monitoramento de performance, análise de riscos, Power BI (DAX, Power Query), Excel.

\textbf{Ferramentas de Gestão:} Monday, Trello, ServiceNow, Spotfire, Pacote Office.

\textbf{Dados e Tecnologia:} SQL, análise de dados, governança de dados, integração entre áreas de negócio e tecnologia.

\textbf{Atendimento e Comunicação:} Onboarding de usuários, treinamento e capacitação, suporte funcional, comunicação consultiva, mediação entre stakeholders, inglês avançado.

\textbf{Idiomas:} Inglês avançado \textbullet{} Espanhol básico \textbullet{} Coreano básico.

\section*{Formação Acadêmica}

\textbf{PUCPR} -- Graduação em Gestão da Tecnologia da Informação \hfill Previsão de conclusão: 06/2026\\
Curitiba, PR

\textbf{Escola Conquer} -- Pós-graduação em Liderança e Gestão em Tecnologia \hfill Conclusão: 2024\\
Curitiba, PR

\textbf{Universidade Tuiuti do Paraná} -- Graduação em Medicina Veterinária \hfill Conclusão: 06/2019\\
Curitiba, PR

\textbf{Qualittas} -- Pós-graduação em Clínica Médica e Cirúrgica de Animais Exóticos e Pets Não Convencionais \hfill Conclusão: 12/2022\\
Curitiba, PR

\end{document}
